\documentclass[a4paper,11pt]{ltjsarticle}

% --- パッケージの読み込み ---
\usepackage{amsmath, amssymb, amsthm, mathrsfs}
\usepackage{luatexja}
\usepackage{tikz}
\usetikzlibrary{cd, shapes, backgrounds, positioning, calc, arrows.meta, fit, patterns, trees}
\usepackage{hyperref}
\usepackage{xcolor}
\usepackage{listings}
\usepackage{tcolorbox}
\usepackage{booktabs}
\usepackage{multicol}

% --- ハイパーリンク設定 ---
\hypersetup{
    colorlinks=true,
    linkcolor=blue!70!black,
    citecolor=blue!70!black,
    urlcolor=blue!70!black
}

% --- 定理環境の定義 ---
\theoremstyle{definition}
\newtheorem{definition}{定義}[section]
\newtheorem{theorem}[definition]{定理}
\newtheorem{lemma}[definition]{補題}
\newtheorem{example}[definition]{例}
\newtheorem{remark}[definition]{注釈}
\newtheorem{problem}[definition]{問題}

% --- マクロ定義 ---
\newcommand{\N}{\mathbb{N}}
\newcommand{\Set}{\mathbf{Set}}
\newcommand{\Cat}{\mathcal{C}}
\newcommand{\Tower}{\mathcal{T}}
\newcommand{\Layer}[1]{\mathrm{layer}_{#1}}
\newcommand{\MinLayer}{\mathrm{minLayer}}
\newcommand{\Hom}{\mathrm{Hom}}
\newcommand{\Id}{\mathrm{id}}

% --- タイトル情報 ---
\title{\textbf{構造塔の圏論的進化：定義の洗練と普遍性} \\ \large --- Lean 4による形式化の軌跡 ---}
\author{
    Lean Code Generation: \textbf{CODEX (OpenAI)} \\
    TeX Explanation \& Typesetting: \textbf{Gemini (Google)}
}
\date{\today}

\begin{document}

\maketitle

\begin{abstract}
本稿では、Lean 4ファイル群（\texttt{CAT1}, \texttt{CAT2\_advanced}, \texttt{CAT2\_revised}, \texttt{CAT2\_complete}）に基づき、「構造塔 (Structure Tower)」の概念がどのように定義され、圏論的な要請（特に普遍射の一意性）によってどのように進化していったかを解説する。
初期の素朴な定義から出発し、関手や自然変換の導入を経て、射の一意性が保証されないという問題に直面する。その解決策として導入された「最小層 (minLayer)」を持つ構造塔 \texttt{StructureTowerWithMin} が、いかにして完全な圏論的振る舞いを獲得するに至ったかを数学的に詳述する。
\end{abstract}

\tableofcontents
\newpage

% =========================================================================
% Section 1: CAT1 - 構造塔の圏
% =========================================================================
\section{構造塔の誕生と圏の構成}
\label{sec:cat1}
\textit{Based on: \texttt{CAT1.md}}

最初のステップとして、集合に「階層構造」を入れた最も基本的な形を定義する。

\subsection{構造塔の定義}
\begin{definition}[Structure Tower]
型 $u, v$ に対し、構造塔 $\Tower$ は以下の要素からなる。
\begin{itemize}
    \item \textbf{基礎集合} (carrier): $X$
    \item \textbf{添字集合} (Index): $I$ （前順序 $\le$ を持つ）
    \item \textbf{層} (layer): 写像 $\Layer{\cdot}: I \to \text{Set } X$
\end{itemize}
これらは以下の公理を満たす。
\begin{enumerate}
    \item \textbf{被覆性}: $\forall x \in X, \exists i \in I, x \in \Layer{i}$
    \item \textbf{単調性}: $i \le j \implies \Layer{i} \subseteq \Layer{j}$
\end{enumerate}
\end{definition}

\subsection{射の定義と圏}
構造塔の間の射 $(f, \phi): \Tower \to \Tower'$ は、以下の対である。
\begin{itemize}
    \item $f: X \to X'$ （基礎集合の写像）
    \item $\phi: I \to I'$ （添字の単調写像）
    \item 条件: $x \in \Layer{i} \implies f(x) \in \Layer{\phi(i)}'$
\end{itemize}
これにより、構造塔全体は圏をなす。

\begin{figure}[h]
\centering
\begin{tikzpicture}
    % Tower T
    \node at (0, 3) {$\Tower$};
    \draw[fill=blue!10] (0,0) ellipse (1.5 and 1);
    \node at (0, -0.5) {$\Layer{i}$};
    \fill (0,0) circle (2pt) node[above] {$x$};

    % Tower T'
    \node at (5, 3) {$\Tower'$};
    \draw[fill=red!10] (5,0) ellipse (2 and 1.5);
    \node at (5, -1) {$\Layer{\phi(i)}'$};
    \fill (5,0.5) circle (2pt) node[above] {$f(x)$};

    % Morphism
    \draw[->, thick] (1.8, 0.5) -- (3.2, 0.5) node[midway, above] {$f$};
    \node at (2.5, -0.5) {構造の保存};
\end{tikzpicture}
\caption{構造塔の射。ある層に含まれる要素は、写像先でも対応する層に含まれなければならない。}
\end{figure}

% =========================================================================
% Section 2: CAT2_advanced - 関手と自然変換
% =========================================================================
\section{圏論的道具立て：関手と自然変換}
\label{sec:cat2_advanced}
\textit{Based on: \texttt{CAT2\_advanced.md}}

圏が定義されたので、より高度な圏論的性質を探る。

\subsection{忘却関手}
構造塔の圏から集合の圏への忘却関手が定義される。
\begin{itemize}
    \item \texttt{forgetCarrier}: $\Tower \mapsto X$ （構造を忘れて基礎集合を見る）
    \item \texttt{forgetIndex}: $\Tower \mapsto I$ （構造を忘れて添字集合を見る）
\end{itemize}

\subsection{自由構造塔 (Free Structure Tower)}
集合 $S$ から生成される「最も緩い」構造塔を考える。
添字集合を $S$ 自身とし、順序を等号のみ（離散順序）とする。
\[ \Layer{s} = \{s\} \]
これは、各要素が自分自身の層にのみ属する構造である。

\begin{problem}[普遍性の予感]
自由構造塔は、忘却関手に対する左随伴になるはずである。つまり、任意の写像 $f: S \to \Tower.\text{carrier}$ に対して、一意的な構造塔の射 $\tilde{f}: \text{Free}(S) \to \Tower$ が存在するだろうか？
\end{problem}

ここで、次のセクションで明らかになる「一意性の欠如」という問題が浮上する。

% =========================================================================
% Section 3: CAT2_revised - 一意性の危機と解決策
% =========================================================================
\section{普遍性の危機と \texttt{StructureTowerWithMin}}
\label{sec:cat2_revised}
\textit{Based on: \texttt{CAT2\_revised.md}}

このファイルは、構造塔の理論における転換点である。初期定義の欠陥を指摘し、修正案を提示している。

\subsection{一意性の欠如}
従来の定義では、要素 $x$ が複数の層 $\Layer{i}, \Layer{j}$ に属することが許されていた。
射 $(f, \phi)$ を構成する際、基礎写像 $f$ が決まっても、添字写像 $\phi$ の取り方には自由度が残ってしまう。

\begin{example}[一意性の破れ]
$\Tower$ において $x \in \Layer{i}$ かつ $x \in \Layer{j}$ ($i \neq j$) とする。
自由構造塔からの射を考えるとき、$x$ に対応する添字を $\phi(x) = i$ とするか $\phi(x) = j$ とするかで、異なる射が定義できてしまう。
これでは、普遍性（$\exists !$）が成り立たない。
\end{example}

\begin{figure}[h]
\centering
\begin{tikzpicture}
    % Ambiguity
    \draw[fill=gray!10] (0,0) circle (2);
    \node at (0, 2.3) {$\Tower$ の層構造};
    
    \draw[blue, thick] (-0.5, 0) circle (1); \node[blue] at (-1, 1) {$\Layer{i}$};
    \draw[red, thick] (0.5, 0) circle (1); \node[red] at (1, 1) {$\Layer{j}$};
    
    \fill (0,0) circle (2pt) node[above] {$x$};
    
    \node[align=left] at (5, 0) {
        $x$ は $i$ にも $j$ にも属する。\\
        射 $\phi$ は $x$ を $i$ に送るべきか？\\
        それとも $j$ に送るべきか？\\
        $\implies$ \textbf{一意に決まらない！}
    };
\end{tikzpicture}
\caption{多重所属による曖昧さ。これが普遍性を妨げる原因となる。}
\end{figure}

\subsection{解決策：最小層 (minLayer) の導入}
この問題を解決するために、\textbf{Version A: StructureTowerWithMin} が提案される。
これは、各要素 $x$ に対して「最小の層」を一意に指定する関数 $\MinLayer$ を構造の一部として組み込むものである。

\begin{definition}[StructureTowerWithMin]
構造塔に加え、以下の関数を要求する。
\[ \MinLayer: X \to I \]
公理：
\begin{enumerate}
    \item $x \in \Layer{\MinLayer(x)}$ （正当性）
    \item $x \in \Layer{i} \implies \MinLayer(x) \le i$ （最小性）
\end{enumerate}
さらに、射 $(f, \phi)$ に対して\textbf{最小層の保存}を要請する：
\[ \MinLayer'(f(x)) = \phi(\MinLayer(x)) \]
\end{definition}

この追加条件により、基礎写像 $f$ が決まれば、添字写像 $\phi$ は上記の方程式によって\textbf{一意に決定される}。これで圏論的な普遍性が回復する。

% =========================================================================
% Section 4: CAT2_complete - 完成された理論
% =========================================================================
\section{完全な形式化：普遍性と直積}
\label{sec:cat2_complete}
\textit{Based on: \texttt{CAT2\_complete.md}}

修正された定義に基づき、完全な圏論的性質を証明する。

\subsection{自由構造塔の普遍性（完全証明）}
半順序集合 $S$ から生成される自由構造塔 $\text{Free}(S)$ を考える。
\begin{theorem}[自由構造塔の普遍性]
任意の構造塔 $\Tower$ と単調写像 $f: S \to \Tower.\text{carrier}$（ただし $\MinLayer$ と整合的）に対し、以下を満たす射 $\varphi: \text{Free}(S) \to \Tower$ が\textbf{ただ一つ}存在する。
\[ \varphi.\text{map}(s) = f(s) \]
\end{theorem}
この定理は、Lean内で `unique` （$\exists !$）を用いて完全に証明されている。

\subsection{直積構造塔と普遍性}
二つの構造塔 $\Tower_1, \Tower_2$ の直積 $\Tower_1 \times \Tower_2$ を定義できる。
\begin{itemize}
    \item 基礎集合：$X_1 \times X_2$
    \item 添字集合：$I_1 \times I_2$
    \item 層：$\Layer{(i,j)} = \Layer{i} \times \Layer{j}$
    \item 最小層：$\MinLayer(x, y) = (\MinLayer(x), \MinLayer(y))$
\end{itemize}

この直積に対しても、射影 $\pi_1, \pi_2$ を伴う普遍性が成立する。すなわち、任意の $\Tower$ からの射 $f_1, f_2$ に対し、$\langle f_1, f_2 \rangle$ が一意に定まる。

\begin{figure}[h]
\centering
\begin{tikzpicture}[node distance=2.5cm, auto]
    \node (T) {$\Tower$};
    \node (P) [right=of T] {$\Tower_1 \times \Tower_2$};
    \node (T1) [above right=of P] {$\Tower_1$};
    \node (T2) [below right=of P] {$\Tower_2$};
    
    \draw[->] (T) to node {$\exists ! \varphi$} (P);
    \draw[->] (P) to node [swap] {$\pi_1$} (T1);
    \draw[->] (P) to node {$\pi_2$} (T2);
    \draw[->, bend left] (T) to node {$f_1$} (T1);
    \draw[->, bend right] (T) to node [swap] {$f_2$} (T2);
\end{tikzpicture}
\caption{直積の普遍性。$\MinLayer$ の保存条件があるおかげで、図式を可換にする射 $\varphi$ は一意に定まる。}
\end{figure}

\section{結論}

一連のファイルは、数学的な構造を形式化する際の「定義の微調整」の重要性を物語っている。
\begin{itemize}
    \item 素朴な集合論的定義 (\texttt{CAT1}) では、一意性が保証されない。
    \item 一意性を保証するために、構造を強化（\texttt{minLayer} の追加）する必要があった (\texttt{CAT2\_revised})。
    \item その結果、美しい圏論的性質（随伴、極限）を持つ理論が完成した (\texttt{CAT2\_complete})。
\end{itemize}
これは、ブルバキ流の構造主義をLeanで実践する過程で得られた、深い数理的洞察と言える。

\end{document}